\section{Design of the Reusable Capital Decision Menu}
\label{app:r16menudesign}
This experiment evaluates the finite Bellman object in
\eqref{eq:r16mainmenu}. A current task changes the consumption utility weight
$w$, the coefficient $\eta$ on squared mean consumption, or the current upper
consumption bound $u$. All future primitives and the feasible future schedule
remain fixed within a calibration. The continuation identity in
Proposition~\ref{prop:r16menucache} therefore holds across tasks. A change to
the future horizon, technology, preferences, innovation law, terminal
dispersion cost, or schedule defines a different continuation and requires a
different cache or fitted object.

\subsection{Calibrations and the query population}
Table~\ref{tab:r16menucalibrations} records all four calibrations. The drift
uses the original fixed dense coupling matrix in dimensions ten and fifty,
with $f(y)=\kappa\tanh(By)$. Productivity is $0.1$, discounting is $0.04$,
the future consumption penalty coefficient is $0.2$, and the original
consumption interval is $[0.02,2]$. The terminal penalty is proportional to
cross-sector dispersion. The period weights, feasible reference schedule,
drift constant, innovation amplitudes, and terminal coefficients are stored
before fitting. Their binary64 values, together with the stored states and
coupling matrices, are the exact real coefficients of the finite economic
problem under study. The continuous-diffusion primitives motivate these
coefficients; they do not make the coarse finite value an unqualified
continuous-time approximation.

\begin{table}[!htbp]
\centering\small
\caption{Complete capital-menu calibration design}
\label{tab:r16menucalibrations}
\begin{tabular}{lrrrrrrr}
\toprule
Calibration & $T$ & $K$ & $h$ & $\sigma_I$ & $\sigma_C$ & $\kappa$ & $\chi$\\
\midrule
Original & 1 & 4 & 0.25 & 0.15 & 0.10 & 0.15 & 0.03\\
Quarterly & 2 & 8 & 0.25 & 0.60 & 0.30 & 0.30 & 0.10\\
Long & 4 & 8 & 0.50 & 0.60 & 0.30 & 0.30 & 0.10\\
Intermediate & 3 & 12 & 0.25 & 0.375 & 0.20 & 0.225 & 0.06\\
\bottomrule
\end{tabular}
\par\medskip\noindent\footnotesize\raggedright $K$ is the number of finite decision
periods and $h=T/K$. Each row is run at $d=10,50$ with sixteen fixed streams.
Original retains the original low-volatility primitive anchor. Quarterly
was used in one disclosed exploratory five-method pilot; Long and
Intermediate had no method-development or outcome pilot. Long increases
calendar time and the length of a period relative to Quarterly while
retaining eight periods.
\end{table}

Nine tasks are the Cartesian product
$w\in\{0.8,1,1.2\}$ and $\eta\in\{0.1,0.2,0.4\}$ with $u=2$.
The other three are $(w,\eta,u)=(0.9,0.3,0.6)$,
$(1.1,0.15,0.8)$, and $(1,0.25,1)$. Every task retains the lower bound
$0.02$. The first nine appeared in the disclosed pilot; the three tighter-cap
tasks did not. Each dimension has thirty-two stored states crossed with all
twelve tasks. The inferential population is the uniform finite collection
of these 384 queries and the sixteen declared streams, separately in each
calibration and dimension. No query, stream, task, or work stage is removed
because of its realized performance. Training inputs, method failures,
completed prefixes, and feasible reference fallbacks remain recorded.

\subsection{Five reusable procedures}
The NBO scalar critic and the vector costate predictor both have width
thirty-two. At each work stage they share the same replay states, uniform
feasible training actions, postdecision arguments, and four-path antithetic
continuation labels. The scalar critic fits gradients of one differentiable
continuation; an intercept fits its level using the value labels. The vector
predictor fits costate components without the restriction that they be the
gradient of a scalar. These procedures consequently share observations while
using different fitted function classes and optimization objectives. Their
economic contrast does not by itself isolate a single statistical mechanism.

The materialized Raw procedure stores costate labels at the reference current
consumption. Its task-conditioned actor maximizes the exact current reward
and the corresponding linearized continuation in the action. Direct policy
optimization also uses a task-conditioned actor, but retains innovations and
recomputes the nonlinear future payoff at each proposed postdecision argument.
Both actors have width thirty-two. Query-level Raw instead optimizes each of
the 384 task-state decisions directly against a sample-average continuation.
Its innovation cache grows from four to sixteen to sixty-four antithetic
paths, and its previous action initializes the next stage. Retaining an
innovation does not eliminate the state-recursion work required by a new
argument. All five methods follow the same cache-validity rule.

For the fitted procedures, the cumulative replay budgets are 128, 512, and
2,048 states. Each field stage has 1,200 Adam updates with learning rate
$0.003$, followed by L-BFGS with at most 300 iterations and 400 function
evaluations. Each actor stage has 1,200 Adam updates with learning rate
$0.003$ and batch size 128. Query optimization uses 180 Adam updates with
learning rate $0.06$. NBO and sample-average Raw then apply L-BFGS to their
scalar objective, with at most sixty iterations and ninety evaluations. The
vector costate procedure retains its 180 gradient-driven action updates;
it has no additional scalar-objective L-BFGS stage. All three completed
stages are assessed. The final payoff bank is inaccessible to fitting and
cannot select a method, checkpoint, work stage, query, or replacement stream.

\subsection{Independent payoff and scalar-accuracy statements}
There are $4\times2\times16=128$ trials, 640 method processes, and 1,920
stage candidates. Conditional on these fitted objects, each task-state query
has 512 independent antithetic confirmation pairs. Common innovations pair
the method contrasts. Proposition~\ref{prop:r16menurange} supplies the
query-conditional range after removing known analytic centers. The numerical
account also retains the fixed Gaussian truncation and arithmetic
allowances. The primary family allocates error $0.018$ across all 216 events:
120 method-versus-reference gains and 96 direct NBO-versus-comparator
contrasts. Thus every comparator has twenty-four direct statements across
all calibrations, dimensions, and stages. There are no additional per-stream
confidence events or continuation-risk estimates.

The economic superiority margin is $10^{-4}$. A protected lower contrast
above zero establishes a positive difference; material superiority requires
the lower contrast to reach $10^{-4}$. Equivalence requires the complete
interval to lie in $[-10^{-4},10^{-4}]$. Positive differences and equivalence
can therefore both be established for the same comparison. These categories
are reported separately. Training loss and an observed economic payoff gain
do not establish a reduction in independent critic mean-squared error.

The separate scalar family allocates error $0.002$ across 128 statements.
For the central task $(w,\eta,u)=(1,0.2,2)$ and the first stored state
(index zero), a successful fit selects its
candidate using the final NBO critic on 257 fixed grid points
in the interval consisting of the reference current consumption plus or
minus $0.1$, intersected with the action box. A failed fit retains the
predeclared feasible reference fallback. All sectors receive this same
scalar action. The independent verifier uses 65,536 antithetic pairs and the
proved derivative and curvature account to cover the entire continuous
interval, including points between the candidate grid. These bounds concern
this restricted finite scalar intervention. They do not bound regret over
all vector actions, future adapted policies, or the continuous diffusion.

\subsection{Complete work and development history}
The final-stage process clock covers startup, imports, input checking, all
three fitting stages, caches, derivative and optimization work, query
deployment, candidate output, and finalization. Earlier-stage clocks combine
the executed prefix with all measured outer overhead and finalization; they
are conservative allocations of this run. They are not observations from
separate online stopping experiments. The independent confirmation process
is a shared workload covering all five methods, all three stages, and the
scalar assessment. Its full bill is reported separately. Charging that bill
to each method is an explicit upper allocation, not a measurement of five
independently executed verification jobs. Timing summaries describe these
runs; the confidence family makes no claim about an unseen runtime or
optimizer distribution.

The development archive retains the single exploratory five-method menu
pilot, its interrupted preliminary attempt, and all twelve fixed-random-feature
explorations, including their adverse results. Their payoff observations are
not confirmation observations. The final source fixes new streams, the
complete task population, every budget, both margins, failure behavior,
power calculations, and all inference code before the registered fitting
and confirmation. The independent fixed-dictionary identity in
Proposition~\ref{prop:r16menuprojection} describes a precise conditional
statistical benefit and its approximation cost. It neither identifies the
nonlinear fitted network with that projection nor makes the observed
five-procedure economic comparison a causal decomposition of critic risk.
